\documentclass[../../../include/open-logic-section]{subfiles}

\begin{document}

\olfileid{sth}{story}{approach}
\olsection{اکٹھا کرن والا مرحلہ وار طریقہ}

ہِیٹھاں ذرا کھول کے دسیا اے کہ اسیں سیٹاں نوں مرحلیاں وچ کیویں ونڈاں گے:
\begin{quote}
	سیٹ \emph{مرحلیاں} وچ بندے نیں۔ ہر مرحلے $S$ لئی کجھ مرحلے $S$ توں
	\emph{پہلاں} آؤندے نیں۔ مرحلے $S$ اُتے، $S$ توں پہلاں دے مرحلیاں وچ بنے
	سیٹاں دے ہر اکٹھ نوں اک سیٹ بنایا جاندا اے۔ مرحلیاں وچ بنے سیٹاں توں
	بناں ہور کوئی سیٹ نہیں۔ \citep[p.~323]{Shoenfield:AST}
\end{quote}
ایہ \emph{سیٹ دے اکٹھا کرن والے مرحلہ وار تصور} دی اک مُڈھلی جھلک اے۔
\olref[sth][][]{part} وچ جیہڑا رسمی سیٹ نظریہ اسیں پیش کراں گے، اوہ
ایس اُتے ٹکیا ہووے گا۔

ایہ گل ہور کھول کے ویکھیے۔ شون فیلڈ دے بیان مطابق ہر مرحلے اُتے
اسیں نویں سیٹ اوہناں {سیٹاں} توں بناندے آں جیہڑے پچھلے مرحلیاں توں
ساڈے کول سن۔ سو شون فیلڈ دے خیال مطابق، پہلے مرحلے، یعنی مرحلے $0$ اُتے،
کوئی \emph{پچھلا} مرحلہ نہیں۔ ایس لئی \emph{ہور وی صاف} اے کہ ساڈے کول
پچھلے مرحلیاں توں کوئی سیٹ نہیں۔\footnote{اسیں کیوں منّیئے کہ پہلا مرحلہ
\emph{موجود اے}؟ \stagesord{} دا حاشیہ \olref[z][story]{sec} وچ ویکھو۔}
سو اسیں صرف اک سیٹ بناندے آں: اوہ سیٹ جیہدے وچ کوئی !!{element}s نہیں،
$\emptyset$۔ مرحلے $1$ اُتے پچھلے مرحلیاں توں ساڈے کول عین اک سیٹ اے،
سو اکّو نواں سیٹ $\{\emptyset\}$ اے۔ مرحلے $2$ اُتے پچھلے مرحلیاں توں
ساڈے کول دو سیٹ نیں تے اسیں دو نویں سیٹ $\{\{\emptyset\}\}$ تے
$\{\emptyset, \{\emptyset\}\}$ بناندے آں۔ مرحلے $3$ اُتے پچھلے مرحلیاں توں
ساڈے کول چار سیٹ نیں، سو اسیں باراں نویں سیٹ بناندے آں\ldots۔
ایس طرح سیٹاں دی اکٹھا کرن والی مرحلہ وار صورت کجھ ایہو جیہی ہووے گی
(عدد مرحلے دسدے نیں):
\begin{center}
	\begin{tikzpicture}[scale=0.6]
	\tikzset{cut_here/.style={densely dotted}}
	\draw (-4,6) -- (-1, 0)-- (2,6); 
	\draw[cut_here] (-5,8)--(-4,6);
	\draw[cut_here] (2,6)--(3,8);
	\draw(-1.5, 1)--(-0.5, 1);
	\draw(-2, 2)--(0, 2);
	\draw(-2.5, 3)--(0.5,3);
	\draw(-3, 4)--(1, 4);
	\draw(-3.5, 5)--(1.5, 5);
	\draw(-4, 6)--(2, 6);
	\node[label] at (0, 0) {\small 0};
	\node[label] at (0.5, 1) {\small 1};
	\node[label] at (1, 2) {\small 2};
	\node[label] at (1.5, 3) {\small 3};
	\node[label] at (2, 4) {\small 4};
	\node[label] at (2.5, 5) {\small 5};
	\node[label] at (3, 6) {\small 6};
	\end{tikzpicture}
\end{center}
سو اسیں ایہ بیان کیوں منّیئے؟

اک وجہ ایہ اے کہ ایہ بیان سوہنا اے تے ورتنا سوکھا اے۔ سبھ توں ظاہر
بیان، یعنی سادے خاصیت توں سیٹ بناؤن دے اصول، دے ڈِگ پَین توں پچھوں
سانوں کوئی چنگا بیان چاہیدا اے۔

پر ایہ بیان \emph{صرف} سوہنا ای نہیں۔ ساڈے کول چنگی وجہ اے ایہ منّن
دی کہ ایس بیان اُتے ٹکیا ہر سیٹ نظریہ \emph{بے تضاد} ہووے گا۔ وجہ ایہ اے۔

سیٹ دے اکٹھا کرن والے مرحلہ وار تصور مطابق اسیں سیٹ مرحلیاں وچ بناندے آں،
تے اوہناں دے !!{element}s ایسی شےواں ہونیاں چاہیدیاں نیں جیہڑیاں
\emph{پہلاں ای} ساڈے کول سن۔ سو کسے وی مرحلے~$S$ لئی اسیں ایہ سیٹ بنا سکدے آں:
\[
	R_S = \Setabs{x}{x \notin x \text{ تے $x$ مرحلے $S$ توں پہلاں موجود سی}}
\]
رسل دا پیراڈوکس ثابت کرن والی دلیل ہُن ایہ ثابت کرے گی کہ $R_S$ آپ مرحلے
$S$ توں پہلاں ساڈے کول نہیں سی۔ تے ایہ تضاد نہیں۔ نالے جے اسیں سیٹ دا
اکٹھا کرن والا مرحلہ وار تصور منّ لئیے، تاں سانوں ایہ \emph{آس} وی نہیں
سی رکھنی چاہیدی کہ اسیں رسل دا سیٹ آپ بنا سکاں گے۔ اوہ اوہناں سارے
سیٹاں دا سیٹ ہووے گا جیہڑے اپنے آپ دے رکن نہیں تے جیہڑے ’’کدی وی ساڈے
کول ہون گے‘‘۔ مختصر گل ایہ اے کہ (جیویں ثابت کیتا جا سکدا اے) رسل دا سیٹ
نہ بنا سکنا ایس بیان مطابق \emph{حیرانی} دی گل نہیں۔ اسیں تے ایہو \emph{اندازہ}
لاندے ساں۔

%In one sense, then, the cumulative-iterative conception of set yields a response to Russell's Paradox which is rather like the \emph{predicativist}'s response. After all, the predicativist said that the collection of all non-self-membered sets$_n$ is a set$_{n+1}$, and this set$_n$ will not be self-membered. (Indeed, most predicativists will treat it as \emph{ungrammatical} to try to ask whether a set is self-membered.)
%
%But there is an important difference between the cumulative-iterative approach and the predicativist approach: the cumulative-iterative approach treats all of our entities as being \emph{of the same kind}. The predicativist will presumably have an empty set$_0$ (i.e., a set$_0$ with no !!{element}s), and an empty set$_1$ (i.e., a set$_1$ with no !!{element}s), and these will be \emph{different entities}. But on the cumulative-iterative approach, there is just one empty set, $\emptyset$, and it will be available at every stage of the hierarchy. 

%Indeed, the Axiom of Extensionality, stated at the start of this chapter, will hold true of the elements in this hierarchy (so that there can be \emph{only one} ``empty'' set). But I do not intend to use the cumulative-iterative conception to justify Extensionality (see \cite{Boolos1971}). Again: I suggest that we should accept Extensionality, just because we are interested in extensional collections. 

\end{document}
